\Subhead{6.1.1}{Апостериорная плотность вероятности. Функция правдоподобия}

Наблюдение $Y$ — доступные данные; сообщение/сигнал $X$ — неизвестная величина или функция, которую требуется определить. Оптимальность всегда относится к заданным модели и потерям, а не означает универсально «лучший приёмник». Формула Байеса
\begin{equation}
 p(x|y)=\frac{p(y|x)p(x)}{p(y)},\qquad
 p(y)=\int p(y|x)p(x)\dd x
\end{equation}
объединяет априорное распределение и правдоподобие. При фиксированном $y$ функция $p(y|x)$ не обязана интегрироваться в единицу по $x$: это не апостериорная плотность.

При гипотезах $H_0$ (шум) и $H_1$ (сигнал плюс шум) достаточным для оптимального бинарного решения статистическим объектом является отношение правдоподобий
\begin{equation}
 \Lambda(y)=\frac{p(y|H_1)}{p(y|H_0)}.
\end{equation}
Если $\bm y=\bm s+\bm n$, $\bm n$ — центрированный гауссов шум с одинаковой при обеих гипотезах невырожденной ковариацией $C$, раскрытие квадратичных форм даёт
\begin{equation}
 \ln\Lambda=\bm s^{\mathsf T}C^{-1}\bm y-
 \tfrac12\bm s^{\mathsf T}C^{-1}\bm s.
\end{equation}
Таким образом, для известного сигнала решение сводится к корреляции с предварительно выбеленным шаблоном. Ковариация задаёт, каким направлениям наблюдения доверять. При неизвестных параметрах требуется либо усреднение по их априорному распределению, либо отдельно обоснованная процедура, например обобщённое отношение правдоподобий; подстановка «истинного» неизвестного параметра не является реализуемым приёмником.
